
\paragraph{Pre-trained model selection.} LogME \citep{you2021logme} scores transferability of a pre-trained model to a target task from its features alone, without fine-tuning. Model Spider \citep{zhang2023modelspider} learns to rank a model zoo efficiently using per-model embeddings. TransferGraph \citep{li2024transfergraph} reframes model selection as link prediction on a model--dataset graph, learning a regressor over graph embeddings; it directly inspired the dataset-similarity view we use in our first approach, though we do not integrate its trained predictor. ArtifactLinker \citep{yu2026artifactlinker} extends this graph-learning framing to a joint artifact graph spanning models, datasets, papers, and code, and pairs ranking with LLM-agent verification; it is the trained backbone our Retrieval and Evaluation agents both build on. These methods differ from ours in two ways. LogME and Model Spider must run every candidate over the target data, so their cost grows with the model zoo and they ignore the evaluation history the community has already recorded; the graph methods exploit that history but only for datasets that already exist as nodes, and they stop at a predicted score. AutoModel Advisor makes an out-of-graph dataset scorable (by summary-embedding retrieval and inductive node insertion), reasons over the graph's measured evidence with an LLM instead of probing each candidate, and treats the resulting ranking as a hypothesis to verify rather than the answer.

\paragraph{LLM- and agent-based discovery.} HuggingGPT \citep{shen2023hugginggpt} uses an LLM to plan and dispatch tasks to Hugging Face models but does not rank or verify candidates against measured performance. AutoML-Agent \citep{trirat2024automlagent} orchestrates a multi-agent LLM pipeline across the full AutoML lifecycle, from data preparation to deployment. In both, the choice of model rests on the LLM's parametric knowledge of model cards and popularity. AutoModel Advisor narrows the scope to model selection but grounds the LLM's choice in graph-conditioned evidence -- measured scores on similar datasets and GNN predictions -- and filters by user constraints before ranking; our coding-agent baseline, which is given the same candidate pool but no graph evidence, approximates this evidence-free setting (see Limitations).

\paragraph{Closed-loop experimentation.} DS-Agent \citep{guo2024dsagent} augments LLM-driven data science with case-based reasoning over past runs. R\&D-Agent \citep{yang2025rdagent} automates iterative research-and-development cycles by feeding evaluation feedback into subsequent iterations. AutoModel Advisor shares this closed-loop spirit but applies it specifically to model \emph{selection}: rather than iterating on a solution's code or pipeline, the loop iterates on the candidate set, feeding per-model measured results and error analysis (Section~\ref{sec:results}) back into retrieval so the revision run confirms, drops, or replaces candidates based on what actually happened on the target data.
